% !TEX program = pdflatex
% Hugging Face Deep RL Course -- Bonus Unit 3.
% Complete original MDX material; compiled with the existing shared preamble.
% Compile from notes/ru/ with latexmk -pdf unitbonus3ru.tex
\PassOptionsToPackage{T2A}{fontenc}
\PassOptionsToPackage{russian}{babel}
\documentclass[11pt,a4paper]{article}

% Shared typography/layout for Hugging Face Deep RL Course notes.
% Compile with pdfLaTeX / latexmk -pdf with notes/ as the working directory.
% e.g. cd notes && latexmk -pdf unit1.tex

\usepackage[T1]{fontenc}
\usepackage[utf8]{inputenc}
\usepackage[english]{babel}
\usepackage{lmodern}
\usepackage{textcomp}
\usepackage{geometry}
\geometry{a4paper,margin=22mm,headheight=14pt,headsep=8mm,footskip=12mm}

\usepackage{amsmath,amssymb,mathtools}
\usepackage{graphicx,caption,needspace,placeins}
\usepackage{array,booktabs,tabularx}
\usepackage{xcolor}
\usepackage{microtype}
\usepackage{enumitem}
\usepackage{listings}
\usepackage[most]{tcolorbox}
\usepackage{titlesec}
\usepackage{fancyhdr}
\usepackage{hyperref}
\usepackage{bookmark}

\definecolor{ink}{HTML}{243041}
\definecolor{accent}{HTML}{38598C}
\definecolor{secondary}{HTML}{60768A}
\definecolor{codebg}{HTML}{F4F6F8}
\definecolor{codeborder}{HTML}{E3E8EE}
\definecolor{notebg}{HTML}{F3F7FC}

\hypersetup{
  colorlinks=true,
  linkcolor=accent,
  urlcolor=accent,
  citecolor=accent,
  pdfborder={0 0 0}
}
\urlstyle{same}

\setlength{\parindent}{0pt}
\setlength{\parskip}{0.38em plus 0.12em minus 0.05em}

\setlist[itemize]{
  topsep=0.35em,
  itemsep=0.13em,
  leftmargin=1.7em
}
\setlist[enumerate]{
  topsep=0.35em,
  itemsep=0.13em,
  leftmargin=1.8em
}

\setcounter{tocdepth}{2}
\setcounter{secnumdepth}{2}

\titleformat{\section}
  {\Large\bfseries\color{ink}}
  {\thesection}{0.7em}{}

\titleformat{\subsection}
  {\large\bfseries\color{accent}}
  {\thesubsection}{0.7em}{}

\titleformat{\subsubsection}
  {\normalsize\bfseries\color{ink}}
  {\thesubsubsection}{0.6em}{}

\titlespacing*{\section}
  {0pt}{2.2ex plus 0.7ex}{0.9ex}
\titlespacing*{\subsection}
  {0pt}{1.7ex plus 0.4ex}{0.6ex}
\titlespacing*{\subsubsection}
  {0pt}{1.3ex plus 0.3ex}{0.45ex}

\pagestyle{fancy}
\fancyhf{}

\fancyhead[L]{
  \footnotesize\sffamily\color{secondary}
  HF Deep RL Course \enspace/\enspace Unit 1
}
\fancyhead[R]{
  \footnotesize\sffamily\color{secondary}
  \nouppercase{\rightmark}
}
\fancyfoot[C]{
  \footnotesize\color{secondary}\thepage
}
\renewcommand{\headrulewidth}{0.25pt}

\captionsetup{
  font=small,
  labelfont={bf,color=secondary},
  textfont={color=ink},
  justification=centering,
  skip=4pt,
  hypcap=false
}

\lstdefinestyle{hfcode}{
  basicstyle=\ttfamily\footnotesize,
  backgroundcolor=\color{codebg},
  frame=single,
  rulecolor=\color{codeborder},
  numbers=none,
  showstringspaces=false,
  breaklines=true,
  breakatwhitespace=false,
  columns=fullflexible,
  keepspaces=true,
  xleftmargin=0pt,
  xrightmargin=0pt,
  aboveskip=0.6em,
  belowskip=0.6em,
  framexleftmargin=7pt,
  framexrightmargin=7pt,
  framextopmargin=5pt,
  framexbottommargin=5pt,
  upquote=true,
  tabsize=4,
}
\lstset{style=hfcode}

\newcommand{\CourseFigure}[2]{%
  \par\addvspace{0.55em}\noindent%
  \begin{minipage}{\linewidth}%
    \centering
    \includegraphics[
      width=0.86\linewidth,
      height=0.295\textheight,
      keepaspectratio
    ]{#1}\par
    \captionof{figure}{#2}%
  \end{minipage}%
  \par\addvspace{0.55em}%
}

\newtcolorbox{coursenote}{
  enhanced,
  breakable,
  colback=notebg,
  colframe=accent!30,
  boxrule=0.35pt,
  arc=1.2mm,
  left=9pt,
  right=9pt,
  top=6pt,
  bottom=6pt
}

\renewenvironment{quote}
  {\begin{coursenote}}
  {\end{coursenote}}

\emergencystretch=2.5em

\providecommand{\tightlist}{
  \setlength{\itemsep}{0pt}
  \setlength{\parskip}{0pt}
}

\providecommand{\passthrough}[1]{#1}

\Urlmuskip=0mu plus 6mu\relax

\renewcommand{\rmdefault}{cmr}
\renewcommand{\sfdefault}{cmss}
\renewcommand{\ttdefault}{cmtt}
\graphicspath{{../assets/unitbonus3/}}
\fancyhead[L]{\footnotesize\sffamily\color{secondary}Курс Hugging Face Deep RL / Бонусный модуль 3}
\fancyhead[R]{\footnotesize\sffamily\color{secondary}Дополнительные темы}
\title{\textbf{Глубокое обучение с подкреплением}\\[0.35em]\large Бонусный модуль 3: дополнительные темы и направления исследований}
\author{Курс Hugging Face Deep RL\\\small Полная русскоязычная учебная версия на основе MDX}
\date{}
\begin{document}
\selectlanguage{russian}
\fontencoding{T2A}\selectfont
\maketitle
\vspace{-0.65em}
\begin{center}\small\color{secondary}Источник: курс Hugging Face Deep RL (оригинальный текст на английском языке).\end{center}
\begingroup\small\setlength{\parskip}{0pt}\tableofcontents\endgroup
\vspace{0.4em}
\section{Введение}\label{introduction}

\CourseFigure{thumbnail.png}{Обзор бонусного модуля 3 курса Hugging Face Deep RL --- источник: \protect\href{https://huggingface.co/datasets/huggingface-deep-rl-course/course-images/resolve/main/en/unit12/thumbnail.png}{[оригинал изображения]}}

Поздравляем с завершением курса! \textbf{Теперь у вас есть прочная база знаний по глубокому обучению с подкреплением}. Но это лишь начало вашего пути: впереди множество интересных направлений. В этом необязательном модуле мы \textbf{собрали материалы для знакомства с различными концепциями и исследовательскими темами в RL}.

В отличие от других модулей, этот раздел подготовили сразу несколько участников команды Hugging Face. Для каждой темы указан её автор.

Звучит интересно? Тогда начнём!

\section{(Автоматическое) обучение по учебному плану в RL}\label{automatic-curriculum-learning-for-rl}

Хотя большинство изученных в этом курсе методов RL хорошо работают на практике, в некоторых случаях их недостаточно. Например, если:

\begin{itemize}
\item
  задача сложная и требует \textbf{постепенного освоения навыков} (скажем, прежде чем научиться преодолевать сложные препятствия, двуногий агент должен научиться стоять, затем ходить и, возможно, прыгать\ldots{});
\item
  среда изменяется, что влияет на сложность задачи, и требуется, чтобы агент был \textbf{устойчив} к таким изменениям.
\end{itemize}

\CourseFigure{bipedal_frame.png}{Обучение двуногого агента ходьбе в TeachMyAgent: пример поэтапного обучения --- источник: \protect\href{https://huggingface.co/datasets/huggingface-deep-rl-course/course-images/resolve/main/en/unit9/bipedal.gif}{[оригинал изображения]}}

\CourseFigure{movable_creepers_frame.png}{Движущиеся препятствия в тестовой среде TeachMyAgent --- источник: \protect\href{https://developmentalsystems.org/TeachMyAgent/}{TeachMyAgent} \protect\href{https://huggingface.co/datasets/huggingface-deep-rl-course/course-images/resolve/main/en/unit9/movable_creepers.gif}{[оригинал изображения]}}

В таких ситуациях агенту нужно последовательно предлагать разные задачи, организованные так, чтобы он постепенно приобретал навыки. Этот подход называется \textbf{Curriculum Learning --- обучением по учебному плану}; как правило, он предполагает вручную составленную последовательность задач. Например, можно изменять генерацию среды или начальные состояния либо использовать Self-Play, управляя сложностью предлагаемых агенту соперников.

Поскольку спроектировать такой учебный план не всегда просто, направление \textbf{Automatic Curriculum Learning (ACL) занимается методами, которые автоматически учатся выбирать и организовывать задачи для повышения эффективности RL-агента}. Портелас и соавторы предложили следующее определение ACL:

\begin{quote}
\ldots{} семейство механизмов, которые автоматически адаптируют распределение обучающих данных, учась подбирать учебные ситуации в соответствии с возможностями RL-агентов.
\end{quote}

Например, OpenAI использовала \textbf{рандомизацию домена (Domain Randomization)} --- случайные изменения параметров среды --- для обучения роботизированной руки собирать кубик Рубика.

\CourseFigure{dr.jpg}{Рандомизация домена при обучении роботизированной руки собирать кубик Рубика --- источник: \protect\href{https://openai.com/blog/solving-rubiks-cube/}{OpenAI: роботизированная рука собирает кубик Рубика} \protect\href{https://huggingface.co/datasets/huggingface-deep-rl-course/course-images/resolve/main/en/unit9/dr.jpg}{[оригинал изображения]}}

Наконец, вы можете сами проверить устойчивость агентов, обученных в \href{https://huggingface.co/spaces/flowers-team/Interactive_DeepRL_Demo}{TeachMyAgent}: менять параметры среды или даже рисовать рельеф.

\CourseFigure{demo.png}{Интерактивная среда TeachMyAgent с изменяемым рельефом --- источник: \protect\href{https://huggingface.co/spaces/flowers-team/Interactive_DeepRL_Demo}{https://huggingface.co/spaces/flowers-team/Interactive\_DeepRL\_Demo} \protect\href{https://huggingface.co/datasets/huggingface-deep-rl-course/course-images/resolve/main/en/unit9/demo.png}{[оригинал изображения]}}

\subsection{Дополнительные материалы}\label{further-reading}

Для более подробного знакомства с темой рекомендуем следующие материалы:

\subsubsection{Обзор направления}\label{overview-of-the-field}

\begin{itemize}
\item
  \href{https://arxiv.org/pdf/2003.04664.pdf}{Автоматическое построение учебного плана для Deep RL: краткий обзор}
\item
  \href{https://lilianweng.github.io/posts/2020-01-29-curriculum-rl/}{Поэтапное обучение с подкреплением}
\end{itemize}

\subsubsection{Современные методы}\label{recent-methods}

\begin{itemize}
\item
  \href{https://arxiv.org/abs/2203.01302}{Эволюция учебных планов через проектирование сред на основе сожаления}
\item
  \href{https://proceedings.mlr.press/v162/klink22a.html}{Обучение по учебному плану с ограниченным оптимальным транспортом}
\item
  \href{https://arxiv.org/abs/2010.03934}{Приоритетное повторное использование уровней}
\end{itemize}

\subsection{Автор раздела}\label{author}

Автор раздела: \href{https://twitter.com/ClementRomac}{Клеман Ромак (Clément Romac)}.

\section{Обобщающая способность в обучении с подкреплением}\label{generalization-in-reinforcement-learning}

Обобщающая способность играет важнейшую роль в обучении с подкреплением. Хотя \textbf{алгоритмы RL могут хорошо работать в контролируемых средах}, реальный мир создаёт \textbf{особые трудности из-за нестационарности и отсутствия заранее заданных границ возможных ситуаций}.

Поэтому для применения RL в реальном мире принципиально важно разрабатывать алгоритмы, устойчивые к изменению условий, способные переносить знания на новые, но похожие задачи и адаптироваться к ранее неизвестным ситуациям.

Если вы хотите глубже изучить эту исследовательскую тему, рекомендуем начать со следующих материалов:

\begin{itemize}
\item
  \href{https://robertkirk.github.io/2022/01/17/generalisation-in-reinforcement-learning-survey.html}{Обобщающая способность в RL: обзор Роберта Кирка}: в этом подробном обзоре даётся \textbf{полезное введение в понятие обобщающей способности в RL}, поэтому он отлично подходит для начала изучения темы.
\item
  \href{https://blog.research.google/2021/09/improving-generalization-in.html?m=1}{Улучшение обобщения в RL с помощью вложений сходства политик}
\end{itemize}

\section{Модельно-ориентированное обучение с подкреплением (MBRL)}\label{model-based-reinforcement-learning-mbrl}

Модельно-ориентированное обучение с подкреплением отличается от подходов без модели прежде всего тем, что агент изучает \emph{модель динамики}. Однако это отличие существенно влияет на дальнейший процесс принятия решений.

Обычно такая модель описывает переходы между состояниями среды: \(s_{t+1} = f_\theta (s_t, a_t)\). Кроме того, в этом подходе можно применять модели обратной динамики (переход от состояния к действию) и модели вознаграждения (прогнозирование награды).

\subsection{Простое определение}\label{simple-definition}

\begin{itemize}
\item
  Агент многократно пытается решить задачу, \textbf{накапливая данные о состояниях и действиях}.
\item
  На основе этих данных агент создаёт средство для рассуждений об окружающем мире --- \emph{модель динамики}.
\item
  С помощью модели динамики агент \textbf{решает, как действовать, прогнозируя будущее}.
\item
  Выполняя действия, \textbf{агент собирает новые данные, улучшает модель и, в идеале, повышает качество последующих действий}.
\end{itemize}

\subsection{Формальное определение}\label{academic-definition}

В модельно-ориентированном обучении с подкреплением (MBRL) агент взаимодействует со средой, \textbf{изучает модель этой среды}, а затем \textbf{использует её для управления (принятия решений)}.

В частности, агент действует в марковском процессе принятия решений (MDP) с функцией перехода \(s_{t+1} = f (s_t , a_t)\) и получает награду на каждом шаге \(r(s_t, a_t)\). На основе собранного набора данных \(D :={ s_i, a_i, s_{i+1}, r_i}\) агент обучает модель \(s_{t+1} = f_\theta (s_t , a_t)\), \textbf{минимизируя отрицательное логарифмическое правдоподобие наблюдаемых переходов}.

Мы используем метод управления с прогнозирующей моделью по выборкам (MPC). Обученная модель динамики позволяет максимизировать ожидаемое вознаграждение на конечном, рекурсивно прогнозируемом горизонте \(\tau\), рассматривая набор действий, выбранных из равномерного распределения \(U(a)\) (см. \href{https://arxiv.org/pdf/2002.04523}{работу}, \href{https://arxiv.org/pdf/2012.09156.pdf}{работу} или \href{https://arxiv.org/pdf/2009.01221.pdf}{работу}).

\subsection{Дополнительные материалы}\label{further-reading-1}

Для более подробного изучения MBRL рекомендуем следующие материалы:

\begin{itemize}
\item
  \href{https://www.natolambert.com/writing/debugging-mbrl}{Публикация в блоге об отладке MBRL}.
\item
  \href{https://arxiv.org/abs/2006.16712}{Обзорная научная статья по MBRL}.
\end{itemize}

\subsection{Автор раздела}\label{author-1}

Автор раздела: \href{https://twitter.com/natolambert}{Нейтан Ламберт (Nathan Lambert)}.

\section{Офлайн- и онлайн-обучение с подкреплением}\label{offline-vs.-online-reinforcement-learning}

Глубокое обучение с подкреплением (RL) --- подход \textbf{к созданию агентов, принимающих решения}. Агент старается изучить оптимальное поведение (политику), взаимодействуя со средой \textbf{методом проб и ошибок и получая вознаграждение как единственный сигнал обратной связи}.

Цель агента --- \textbf{максимизировать суммарное вознаграждение}, называемое возвратом (return). В основе RL лежит \emph{гипотеза о вознаграждении}: любую цель можно описать как \textbf{максимизацию математического ожидания суммарного вознаграждения}.

Агенты глубокого обучения с подкреплением \textbf{обучаются на пакетах опыта}. Но откуда берутся эти данные?

\CourseFigure{offlinevsonlinerl_frame.png}{Сравнение офлайн- и онлайн-обучения с подкреплением --- источник: \protect\href{https://offline-rl.github.io/}{этот материал} \protect\href{https://huggingface.co/datasets/huggingface-deep-rl-course/course-images/resolve/main/en/unit12/offlinevsonlinerl.gif}{[оригинал изображения]}}

\begin{itemize}
\tightlist
\item
  При \emph{онлайн-обучении с подкреплением}, которое мы изучали на протяжении курса, агент \textbf{непосредственно собирает данные}: получает пакет опыта, \textbf{взаимодействуя со средой}, а затем сразу использует его для обновления политики (или сохраняет опыт в буфере воспроизведения для дальнейшего обучения).
\end{itemize}

Однако это означает, что нужно либо \textbf{обучать агента непосредственно в реальном мире, либо иметь симулятор}. Если симулятора нет, его придётся построить. Это может быть сложно (как достоверно воспроизвести всю сложность реального мира?), дорого и небезопасно: агент способен использовать ошибки симулятора для получения преимущества.

\begin{itemize}
\tightlist
\item
  При \emph{офлайн-обучении с подкреплением}, напротив, агент \textbf{использует только данные, собранные другими агентами или полученные из демонстраций людей}. Во время обучения он \textbf{не взаимодействует со средой}.
\end{itemize}

Процесс выглядит следующим образом:

\begin{itemize}
\item
  \textbf{Создайте набор данных} с помощью одной или нескольких политик и/или действий человека.
\item
  \textbf{Обучите политику методом офлайн-RL}, используя этот набор данных.
\end{itemize}

У этого подхода есть недостаток: \emph{проблема контрфактических запросов}. Что делать, если агент \textbf{решит выполнить действие, для которого в наборе данных нет примеров}? Например, он захочет повернуть направо на перекрёстке, но соответствующая траектория отсутствует.

Для этой проблемы существуют решения. Если хотите глубже изучить офлайн-обучение с подкреплением, можете \href{https://www.youtube.com/watch?v=k08N5a0gG0A}{посмотреть это видео}.

\subsection{Дополнительные материалы}\label{further-reading-2}

Для более подробного знакомства с темой рекомендуем следующие материалы:

\begin{itemize}
\item
  \href{https://www.youtube.com/watch?v=qgZPZREor5I}{Офлайн-обучение с подкреплением: лекция Сергея Левина}
\item
  \href{https://arxiv.org/abs/2005.01643}{Офлайн-обучение с подкреплением: руководство, обзор и открытые проблемы}
\end{itemize}

\subsection{Автор раздела}\label{author-2}

Автор раздела: \href{https://twitter.com/ThomasSimonini}{Томас Симонини (Thomas Simonini)}.

\section{Трансформеры принятия решений (Decision Transformers)}\label{decision-transformers}

Модель Decision Transformer представлена в работе \href{https://arxiv.org/abs/2106.01345}{«Decision Transformer: Reinforcement Learning via Sequence Modeling» Чена и соавторов}. Она рассматривает обучение с подкреплением как задачу условного моделирования последовательностей.

Основная идея заключается в следующем. Вместо того чтобы обучать политику методами RL (например, аппроксимировать функцию ценности, подсказывающую, какое действие максимизирует возврат), \textbf{мы используем алгоритм моделирования последовательностей --- Transformer}. Получив желаемое значение возврата, предыдущие состояния и действия, он генерирует будущие действия, направленные на достижение этого возврата. Это авторегрессионная модель, обусловленная целевым возвратом, предшествующими состояниями и действиями.

Это существенное изменение парадигмы RL: вместо традиционных алгоритмов используется генеративное моделирование траекторий, то есть совместного распределения последовательностей состояний, действий и наград. Таким образом, Decision Transformer не максимизирует возврат напрямую, а генерирует последовательность будущих действий, которая должна привести к заданному возврату.

Команда Hugging Face Transformers интегрировала Decision Transformer --- метод офлайн-RL --- в библиотеку Transformers и Hugging Face Hub.

\subsection{Как устроены трансформеры принятия решений}\label{learn-about-decision-transformers}

Чтобы лучше понять Decision Transformer, прочитайте наш материал \href{https://huggingface.co/blog/decision-transformers}{«Знакомство с Decision Transformers на Hugging Face»}.

\subsection{Обучите свой первый Decision Transformer}\label{train-your-first-decision-transformers}

После знакомства с принципом работы Decision Transformer по \href{https://huggingface.co/blog/decision-transformers}{нашему вводному материалу} вы готовы самостоятельно обучить первую модель на офлайн-данных, чтобы она научила виртуального гепарда (HalfCheetah) бежать.

Начните с руководства: \url{https://huggingface.co/blog/train-decision-transformers}

\subsection{Дополнительные материалы}\label{further-reading-3}

Для дальнейшего изучения рекомендуем следующие материалы:

\begin{itemize}
\item
  \href{https://arxiv.org/abs/2106.01345}{Decision Transformer: RL через моделирование последовательностей}
\item
  \href{https://arxiv.org/abs/2202.05607}{Онлайн-вариант Decision Transformer}
\end{itemize}

\subsection{Автор раздела}\label{author-3}

Автор раздела: \href{https://twitter.com/edwardbeeching}{Эдвард Бичинг (Edward Beeching)}.

\section{Обучение с подкреплением на основе обратной связи от человека (RLHF)}\label{rlhf}

Обучение с подкреплением на основе обратной связи от человека (RLHF) --- это \textbf{метод интеграции человеческих оценок в процесс оптимизации на основе RL}. Его применение связано с \textbf{трудностью формального моделирования человеческих предпочтений}.

Для многих задач даже при попытке записать формулу идеального результата предпочтения разных людей будут отличаться.

Обновление моделей \textbf{на основе измеряемых данных позволяет частично смягчить эти проблемы машинного обучения, неразрывно связанные с человеческими предпочтениями}.

\subsection{С чего начать изучение RLHF}\label{start-learning-about-rlhf}

Чтобы приступить к изучению RLHF:

\begin{enumerate}
\def\labelenumi{\arabic{enumi}.}
\item
  Прочитайте вводную статью: \href{https://huggingface.co/blog/rlhf}{«Иллюстрированное введение в обучение с подкреплением по обратной связи от человека (RLHF)»}.
\item
  Посмотрите запись прямого эфира, в котором Нейтан объяснил основы RLHF и то, как эту технологию используют в современных инструментах машинного обучения, таких как ChatGPT. Большая часть лекции посвящена обзору связанных друг с другом ML-моделей. Рассматриваются основы обработки естественного языка, RL и применение RLHF к большим языковым моделям. В завершение обсуждаются открытые вопросы в RLHF.
\end{enumerate}

\textbf{Видеолекция:} \href{https://www.youtube.com/watch?v=2MBJOuVq380}{смотреть на YouTube}

\begin{enumerate}
\def\labelenumi{\arabic{enumi}.}
\setcounter{enumi}{2}
\tightlist
\item
  Прочитайте и другие статьи по теме, например \href{https://robotic.substack.com/p/rlhf-chatgpt-data-moats}{«Продолжение дискуссии о закрытых API и открытом коде: RLHF, ChatGPT и преимущества данных»}. Если знаете интересные публикации, расскажите нам о них!
\end{enumerate}

\subsection{Дополнительная литература}\label{additional-readings}

\emph{Примечание: этот текст взят из упомянутой выше статьи об RLHF}. Ниже перечислены основные научные публикации по RLHF на момент написания материала. Направление стало широко известно с развитием Deep RL (примерно с 2017 года), а затем расширилось в связи с применением больших языковых моделей в крупных технологических компаниях. Начнём с исследований, предшествовавших акценту на языковых моделях:

\begin{itemize}
\item
  \href{https://www.cs.utexas.edu/~pstone/Papers/bib2html-links/ICDL08-knox.pdf}{TAMER: ручное обучение агента с помощью оценочного подкрепления} (Knox и Stone, 2008): предложен агент, для которого люди последовательно оценивают выполненные действия, позволяя обучить модель вознаграждения.
\item
  \href{http://proceedings.mlr.press/v70/macglashan17a/macglashan17a.pdf}{Интерактивное обучение на зависящей от политики обратной связи человека} (MacGlashan и соавт., 2017): предложен актор-критик COACH, в котором положительные и отрицательные оценки людей используются для настройки функции преимущества.
\item
  \href{https://proceedings.neurips.cc/paper/2017/hash/d5e2c0adad503c91f91df240d0cd4e49-Abstract.html}{Глубокое обучение с подкреплением на основе человеческих предпочтений} (Christiano и соавт., 2017): RLHF применяется к предпочтениям между траекториями в играх Atari.
\item
  \href{https://ojs.aaai.org/index.php/AAAI/article/view/11485}{Deep TAMER: интерактивная настройка поведения агента в высокоразмерных пространствах состояний} (Warnell и соавт., 2018): развитие TAMER, использующее глубокую нейросеть для прогнозирования вознаграждения.
\end{itemize}

Ниже --- подборка исследований, демонстрирующих эффективность RLHF для языковых моделей:

\begin{itemize}
\item
  \href{https://arxiv.org/abs/1909.08593}{Дообучение языковых моделей на основе человеческих предпочтений} (Zieglar и соавт., 2019): раннее исследование влияния обучения вознаграждению на четыре конкретные задачи.
\item
  \href{https://proceedings.neurips.cc/paper/2020/hash/1f89885d556929e98d3ef9b86448f951-Abstract.html}{Обучение суммаризации текста на основе обратной связи человека} (Stiennon и соавт., 2020): применение RLHF к краткому изложению текстов. Продолжение --- \href{https://arxiv.org/abs/2109.10862}{Рекурсивное реферирование книг с обратной связью от человека} (команда OpenAI Alignment, 2021).
\item
  \href{https://arxiv.org/abs/2112.09332}{WebGPT: ответы на вопросы с помощью браузера и обратной связи человека} (OpenAI, 2021): RLHF для обучения агента навигации по веб-страницам.
\item
  InstructGPT: \href{https://arxiv.org/abs/2203.02155}{Обучение языковых моделей следованию инструкциям с помощью обратной связи человека} (команда OpenAI Alignment, 2022): применение RLHF к языковой модели общего назначения (\href{https://openai.com/blog/instruction-following/}{публикация в блоге} об InstructGPT).
\item
  GopherCite: \href{https://www.deepmind.com/publications/gophercite-teaching-language-models-to-support-answers-with-verified-quotes}{Обучение языковых моделей подкреплять ответы проверенными цитатами} (Menick и соавт., 2022): RLHF используется, чтобы языковая модель давала ответы с конкретными ссылками на источники.
\item
  Sparrow: \href{https://arxiv.org/abs/2209.14375}{Улучшение согласования диалоговых агентов с помощью целевых оценок людей} (Glaese и соавт., 2022): дообучение диалогового агента с использованием RLHF.
\item
  \href{https://openai.com/blog/chatgpt/}{ChatGPT: оптимизация языковых моделей для диалога} (OpenAI, 2022): обучение универсального разговорного помощника с использованием RLHF.
\item
  \href{https://arxiv.org/abs/2210.10760}{Законы масштабирования для чрезмерной оптимизации модели вознаграждения} (Gao и соавт., 2022): исследуются закономерности масштабирования обученной модели предпочтений в RLHF.
\item
  \href{https://arxiv.org/abs/2204.05862}{Обучение полезного и безопасного помощника с использованием RLHF} (Anthropic, 2022): подробное описание обучения языковой модели-помощника с обратной связью человека.
\item
  \href{https://arxiv.org/abs/2209.07858}{Красная команда против языковых моделей для снижения вреда: методы, масштабирование и выводы} (Ganguli и соавт., 2022): подробное описание усилий по выявлению, измерению и снижению потенциально вредных ответов языковых моделей.
\item
  \href{https://arxiv.org/abs/2208.02294}{Динамическое планирование в открытом диалоге с помощью обучения с подкреплением} (Cohen и соавт., 2022): RL используется для улучшения разговорных навыков диалогового агента.
\item
  \href{https://arxiv.org/abs/2210.01241}{(Не) подходит ли обучение с подкреплением для обработки естественного языка? Тестовые задачи, базовые методы и компоненты оптимизации языковой политики} (Ramamurthy, Ammanabrolu и соавт., 2022): обсуждается пространство решений при создании открытых инструментов RLHF и предлагается алгоритм NLPO (Natural Language Policy Optimization) как альтернатива PPO.
\end{itemize}

\subsection{Автор раздела}\label{author-4}

Автор раздела: \href{https://twitter.com/natolambert}{Нейтан Ламберт (Nathan Lambert)}.

\section{Языковые модели в RL}\label{language-models-in-rl}

\subsection{Языковые модели содержат полезные знания для агентов}\label{lms-encode-useful-knowledge-for-agents}

\textbf{Языковые модели} (LM) демонстрируют впечатляющие способности при работе с текстом: от ответов на вопросы до пошаговых рассуждений. Кроме того, обучение на огромных текстовых корпусах позволяет им \textbf{сохранять знания разных типов, в том числе абстрактные представления о физических законах нашего мира} (например, что можно сделать с предметом и что произойдёт при его вращении).

Возникает естественный вопрос, который недавно начали активно изучать: могут ли такие знания помочь агентам, в том числе роботам, в выполнении повседневных задач? Хотя работы показали интересные результаты, предложенные агенты не обладали механизмом обучения. \textbf{Из-за этого они не могли адаптироваться к среде (например, исправлять неверные знания) или осваивать новые навыки}.

\CourseFigure{language.png}{Применение знаний языковой модели к возможностям роботизированных систем --- источник: \protect\href{https://ai.googleblog.com/2022/08/towards-helpful-robots-grounding.html}{К полезным роботам: связь языка и возможностей роботов} \protect\href{https://huggingface.co/datasets/huggingface-deep-rl-course/course-images/resolve/main/en/unit12/language.png}{[оригинал изображения]}}

\subsection{Языковые модели и обучение с подкреплением}\label{lms-and-rl}

Поэтому между языковыми моделями, предоставляющими знания о мире, и RL, которое позволяет согласовывать и корректировать эти знания через взаимодействие со средой, возможна полезная синергия. Это особенно важно для RL, поскольку в данной области часто используют постановку \textbf{tabula rasa}: агент изучает всё с нуля, что приводит к:

\begin{enumerate}
\def\labelenumi{\arabic{enumi})}
\item
  Неэффективному использованию обучающих примеров.
\item
  Поведению, которое кажется человеку неожиданным.
\end{enumerate}

Одной из первых попыток стала работа \href{https://arxiv.org/abs/2302.02662v1}{«Привязка больших языковых моделей к среде посредством онлайн-обучения с подкреплением»}. Авторы решали задачу \textbf{адаптации и согласования языковой модели с текстовой средой при помощи PPO}. Было показано, что знания, уже содержащиеся в языковой модели, ускоряют адаптацию к среде (что открывает путь к RL с меньшими затратами выборки) и после согласования помогают лучше обобщать знания на новые задачи.

\textbf{Видеодемонстрация:} \href{https://huggingface.co/datasets/huggingface-deep-rl-course/course-images/resolve/main/en/unit12/papier_v4.mp4}{открыть оригинальное видео}

Другое направление рассматривается в работе \href{https://arxiv.org/abs/2302.06692}{«Использование больших языковых моделей для направления предварительного обучения в RL»}. Языковая модель остаётся замороженной, но её знания применяются, чтобы \textbf{направлять исследование среды RL-агентом}. Это помогает агенту выбирать осмысленные для человека и потенциально полезные действия без необходимости постоянно привлекать человека во время обучения.

\CourseFigure{language2.png}{Использование знаний языковой модели для исследования среды RL-агентом --- источник: \protect\href{https://ai.googleblog.com/2022/08/towards-helpful-robots-grounding.html}{К полезным роботам: связь языка и возможностей роботов} \protect\href{https://huggingface.co/datasets/huggingface-deep-rl-course/course-images/resolve/main/en/unit12/language2.png}{[оригинал изображения]}}

Пока эти результаты предварительны и имеют ряд ограничений: например, наблюдения агента приходится преобразовывать в текст перед передачей языковой модели, а взаимодействие с очень большими языковыми моделями требует значительных вычислительных ресурсов.

\subsection{Дополнительные материалы}\label{further-reading-4}

Для более подробного знакомства с темой рекомендуем следующие материалы:

\begin{itemize}
\item
  \href{https://ai.googleblog.com/2023/02/google-research-2022-beyond-robotics.html}{Google Research: робототехника в 2022 году и далее}
\item
  \href{https://arxiv.org/abs/2202.01771}{Предобученные языковые модели для интерактивного принятия решений}
\item
  \href{https://arxiv.org/abs/2302.02662v1}{Привязка больших языковых моделей к среде через онлайн-RL}
\item
  \href{https://arxiv.org/abs/2302.06692}{Направление предварительного обучения в RL с помощью языковых моделей}
\end{itemize}

\subsection{Автор раздела}\label{author-5}

Автор раздела: \href{https://twitter.com/ClementRomac}{Клеман Ромак (Clément Romac)}.

\section{Знакомство с Unreal Learning Agents}\label{an-introduction-to-unreal-learning-agents}

\href{https://dev.epicgames.com/community/learning/tutorials/8OWY/unreal-engine-learning-agents-introduction}{Learning Agents} --- плагин Unreal Engine (UE), позволяющий \textbf{обучать игровых персонажей на основе методов машинного обучения (ML) непосредственно в Unreal}.

Этот новый интересный плагин позволяет создавать собственные среды в Unreal Engine и обучать в них агентов.

Рассмотрим, как \textbf{начать работу и обучить виртуальный автомобиль управлению движением в среде Unreal Engine}.

\CourseFigure{learning-agents-car.png}{Unreal Engine Learning Agents: обучение агента управлению автомобилем --- источник: \protect\href{https://dev.epicgames.com/community/learning/tutorials/qj2O/unreal-engine-learning-to-drive}{Руководство по обучению автомобиля в Learning Agents} \protect\href{https://huggingface.co/datasets/huggingface-deep-rl-course/course-images/resolve/main/en/unit12/learning-agents-car.png}{[оригинал изображения]}}

\subsection{Случай 1: для начинающих и незнакомых с Unreal Engine}\label{case-1-i-dont-know-anything-about-unreal-engine-and-beginners-in-unreal-engine}

Если вы ещё не знакомы с Unreal Engine, не переживайте! Чтобы начать использовать Learning Agents, предлагаем пройти два вводных курса:

\begin{enumerate}
\def\labelenumi{\arabic{enumi}.}
\item
  Освойте основы: начните с курса \href{https://dev.epicgames.com/community/learning/courses/ZpX/your-first-hour-in-unreal-engine-5/E7L/introduction-to-your-first-hour-in-unreal-engine-5}{«Ваш первый час в Unreal Engine 5»}. Он \textbf{даст базовые знания, необходимые для работы с Unreal}.
\item
  Познакомьтесь с Blueprints: изучите визуальную систему скриптов Unreal Engine. \href{https://youtu.be/W0brCeJNMqk?si=zy4t4t1l6FMIzbpz}{Этот видеокурс} познакомит вас с данным инструментом.
\end{enumerate}

Овладев основами, \textbf{вы готовы экспериментировать с Learning Agents}:

\begin{enumerate}
\def\labelenumi{\arabic{enumi}.}
\setcounter{enumi}{2}
\item
  Получите общее представление о Learning Agents, \href{https://dev.epicgames.com/community/learning/tutorials/8OWY/unreal-engine-learning-agents-introduction}{прочитав этот обзор}.
\item
  \href{https://dev.epicgames.com/community/learning/tutorials/qj2O/unreal-engine-learning-to-drive}{Научите автомобиль ездить с помощью обучения с подкреплением в Learning Agents}.
\item
  \href{https://www.youtube.com/watch?v=NwYUNlFvajQ}{Посмотрите руководство по имитационному обучению с плагином Learning Agents для Unreal Engine 5.3}.
\end{enumerate}

\subsection{Случай 2: для знакомых с Unreal Engine}\label{case-2-im-familiar-with-unreal}

Если вы уже знакомы с Unreal Engine, можете сразу перейти к Learning Agents, воспользовавшись следующими руководствами:

\begin{enumerate}
\def\labelenumi{\arabic{enumi}.}
\item
  Получите общее представление о Learning Agents, \href{https://dev.epicgames.com/community/learning/tutorials/8OWY/unreal-engine-learning-agents-introduction}{прочитав этот обзор}.
\item
  \href{https://dev.epicgames.com/community/learning/tutorials/qj2O/unreal-engine-learning-to-drive}{Научите автомобиль ездить с помощью обучения с подкреплением в Learning Agents}.
\item
  \href{https://www.youtube.com/watch?v=NwYUNlFvajQ}{Посмотрите руководство по имитационному обучению с плагином Learning Agents для Unreal Engine 5.3}.
\end{enumerate}

\section{Godot RL Agents}\label{godot-rl-agents}

\href{https://github.com/edbeeching/godot_rl_agents}{Godot RL Agents} --- пакет с открытым исходным кодом, который даёт разработчикам игр, исследователям ИИ и энтузиастам возможность \textbf{обучать сложному поведению неигровых персонажей (NPC) и агентов}.

Библиотека предоставляет:

\begin{itemize}
\item
  Интерфейс между играми, созданными в \href{https://godotengine.org/}{Godot Engine}, и алгоритмами машинного обучения на Python.
\item
  Обёртки для четырёх известных RL-фреймворков: \href{https://stable-baselines3.readthedocs.io/en/master/}{StableBaselines3}, \href{https://docs.cleanrl.dev/}{CleanRL}, \href{https://www.samplefactory.dev/}{Sample Factory} и \href{https://docs.ray.io/en/latest/rllib-algorithms.html}{Ray RLlib}.
\item
  Поддержку агентов с памятью на основе LSTM или интерфейсов внимания.
\item
  Поддержку \emph{двумерных и трёхмерных игр}.
\item
  Набор \emph{датчиков ИИ}, расширяющих возможности агента по наблюдению игрового мира.
\item
  Godot и Godot RL Agents \textbf{полностью бесплатны, имеют открытый код и распространяются по весьма свободной лицензии MIT}. Без скрытых условий и лицензионных отчислений.
\end{itemize}

Узнать больше о Godot RL Agents можно на \href{https://github.com/edbeeching/godot_rl_agents}{странице GitHub} или из \href{https://arxiv.org/abs/2112.03636}{статьи, представленной на семинаре AAAI 2022}. Создатель библиотеки \href{https://edbeeching.github.io/}{Эд Бичинг} работает исследователем в Hugging Face.

Установка библиотеки проста: \texttt{pip\ install\ godot-rl}

\subsection{Создание собственной RL-среды с Godot RL Agents}\label{create-a-custom-rl-environment-with-godot-rl-agents}

В этом разделе вы \textbf{научитесь создавать собственную среду в игровом движке Godot}, а затем реализуете контроллер ИИ, который учится играть с помощью глубокого обучения с подкреплением.

Сегодня мы создадим простую игру, \textbf{которая тем не менее демонстрирует многие возможности Godot Engine и библиотеки Godot RL Agents}. Затем вы сможете изучить примеры более сложных сред и способов поведения.

Наша среда называется Ring Pong. Это разновидность пинг-понга, где игровое поле имеет форму кольца, а ракетка перемещается по его окружности. \textbf{Цель --- удерживать мяч внутри кольца, отбивая его ракеткой}.

\CourseFigure{ringpong_frame.png}{Ring Pong: пользовательская игровая среда на Godot --- источник: \protect\href{https://huggingface.co/datasets/huggingface-deep-rl-course/course-images/resolve/main/en/unit9/ringpong.gif}{[оригинал изображения]}}

\subsubsection{Установка игрового движка Godot}\label{installing-the-godot-game-engine}

\href{https://godotengine.org/}{Игровой движок Godot} --- инструмент с открытым исходным кодом для \textbf{разработки видеоигр, программ и пользовательских интерфейсов}.

Godot Engine --- многофункциональный кроссплатформенный игровой движок для создания 2D- и 3D-игр в единой среде. Он предоставляет большой набор стандартных инструментов, поэтому разработчики \textbf{могут сосредоточиться на создании игр, не изобретая заново базовые механизмы}. Игры можно экспортировать практически одним нажатием для основных настольных ОС (Linux, macOS, Windows), мобильных платформ (Android, iOS) и веба (HTML5).

Мы проведём вас через реализацию агента, но, возможно, вам захочется лучше изучить сам Godot. Его \href{https://docs.godotengine.org/en/latest/index.html}{документация} весьма подробна, а на YouTube есть множество уроков. Также рекомендуем источники \href{https://www.gdquest.com/}{GDQuest}, \href{https://kidscancode.org/godot_recipes/4.x/}{KidsCanCode} и \href{https://www.youtube.com/channel/UCczi7Aq_dTKrQPF5ZV5J3gg}{Bramwell}.

Чтобы создавать игры в Godot, \textbf{сначала необходимо скачать редактор}. В материалах курса заявлена поддержка версии Godot 4.0 библиотекой Godot RL Agents.

Редактор можно скачать по следующим ссылкам:

\begin{itemize}
\item
  \href{https://downloads.tuxfamily.org/godotengine/4.0.1/Godot_v4.0.1-stable_win64.exe.zip}{Windows}
\item
  \href{https://downloads.tuxfamily.org/godotengine/4.0.1/Godot_v4.0.1-stable_macos.universal.zip}{Mac}
\item
  \href{https://downloads.tuxfamily.org/godotengine/4.0.1/Godot_v4.0.1-stable_linux.x86_64.zip}{Linux}
\end{itemize}

\subsubsection{Загрузка стартового проекта}\label{loading-the-starter-project}

Мы предоставляем два варианта исходного кода:

\begin{itemize}
\item
  \href{https://drive.google.com/file/d/1C7xd3TibJHlxFEJPBgBLpksgxrFZ3D8e/view?usp=share_link}{Стартовый проект: скачайте его и выполняйте шаги руководства}.
\item
  \href{https://drive.google.com/file/d/1k-b2Bu7uIA6poApbouX4c3sq98xqogpZ/view?usp=share_link}{Готовая версия проекта для сравнения и поиска ошибок}.
\end{itemize}

Чтобы загрузить проект, в Godot Project Manager нажмите \textbf{Import}, перейдите в каталог с файлами и выберите \textbf{project.godot}.

Если нажать F5 или запустить игру из редактора, можно играть в режиме человека. Одновременно запускается несколько копий игры, поскольку мы хотим ускорить обучение ИИ-агента с помощью параллельных сред.

\subsubsection{Установка плагина Godot RL Agents}\label{installing-the-godot-rl-agents-plugin}

Плагин Godot RL Agents можно установить из репозитория GitHub или через встроенный каталог Godot AssetLib.

Сначала откройте AssetLib и выполните поиск по слову ``rl.

\CourseFigure{godot1.png}{Поиск плагина RL Agents в каталоге ресурсов Godot AssetLib --- источник: \protect\href{https://huggingface.co/datasets/huggingface-deep-rl-course/course-images/resolve/main/en/unit9/godot1.png}{[оригинал изображения]}}

Затем выберите Godot RL Agents, нажмите Download, снимите галочки с файлов LICENSE и README .md и нажмите Install.

\CourseFigure{godot2.png}{Установка плагина RL Agents в Godot --- источник: \protect\href{https://huggingface.co/datasets/huggingface-deep-rl-course/course-images/resolve/main/en/unit9/godot2.png}{[оригинал изображения]}}

Теперь плагин Godot RL Agents загружен на компьютер. Откройте Project -\textgreater{} Project settings и включите дополнение:

\CourseFigure{godot3.png}{Включение расширения RL Agents в настройках проекта Godot --- источник: \protect\href{https://huggingface.co/datasets/huggingface-deep-rl-course/course-images/resolve/main/en/unit9/godot3.png}{[оригинал изображения]}}

\subsubsection{Добавление контроллера ИИ}\label{adding-the-ai-controller}

Теперь добавим в нашу игру контроллер ИИ. Откройте сцену player.tscn: слева должно появиться дерево узлов примерно следующего вида:

\CourseFigure{godot4.png}{Дерево сцены Godot перед добавлением контроллера ИИ --- источник: \protect\href{https://huggingface.co/datasets/huggingface-deep-rl-course/course-images/resolve/main/en/unit9/godot4.png}{[оригинал изображения]}}

Щёлкните правой кнопкой мыши по узлу \textbf{Player} и выберите \textbf{Add Child Node}. В списке узлов найдите AIController3D и создайте его.

\CourseFigure{godot5.png}{Добавление дочернего узла AIController3D в Godot --- источник: \protect\href{https://huggingface.co/datasets/huggingface-deep-rl-course/course-images/resolve/main/en/unit9/godot5.png}{[оригинал изображения]}}

Узел AIController появится в дереве сцены. Рядом с ним будет значок скрипта: нажмите его, чтобы открыть связанный с контроллером код. Godot использует язык GDScript, синтаксически похожий на Python. Скрипт содержит методы, которые нужно реализовать для работы контроллера ИИ.

\begin{lstlisting}
#-- Methods that need implementing using the "extend script" option in Godot --#
func get_obs() -> Dictionary:
	assert(false, "the get_obs method is not implemented when extending from ai_controller")
	return {"obs":[]}

func get_reward() -> float:
	assert(false, "the get_reward method is not implemented when extending from ai_controller")
	return 0.0

func get_action_space() -> Dictionary:
	assert(false, "the get get_action_space method is not implemented when extending from ai_controller")
	return {
		"example_actions_continous" : {
			"size": 2,
			"action_type": "continuous"
		},
		"example_actions_discrete" : {
			"size": 2,
			"action_type": "discrete"
		},
		}

func set_action(action) -> void:
	assert(false, "the get set_action method is not implemented when extending from ai_controller")
# -----------------------------------------------------------------------------#
\end{lstlisting}

Для реализации этих методов создадим класс, наследующий AIController3D. В Godot это делается легко и называется расширением класса (extending).

Щёлкните правой кнопкой мыши по узлу AIController3D, выберите ``Extend Script и назовите новый скрипт \texttt{controller.gd}. В результате появится почти пустой файл следующего вида:

\begin{lstlisting}
extends AIController3D

# Called when the node enters the scene tree for the first time.
func _ready():
	pass # Replace with function body.

# Called every frame. 'delta' is the elapsed time since the previous frame.
func _process(delta):
	pass
\end{lstlisting}

Теперь реализуем четыре отсутствующих метода: удалим этот шаблон и заменим его следующим кодом:

\begin{lstlisting}
extends AIController3D

# Stores the action sampled for the agent's policy, running in python
var move_action : float = 0.0

func get_obs() -> Dictionary:
	# get the balls position and velocity in the paddle's frame of reference
	var ball_pos = to_local(_player.ball.global_position)
	var ball_vel = to_local(_player.ball.linear_velocity)
	var obs = [ball_pos.x, ball_pos.z, ball_vel.x/10.0, ball_vel.z/10.0]

	return {"obs":obs}

func get_reward() -> float:
	return reward

func get_action_space() -> Dictionary:
	return {
		"move_action" : {
			"size": 1,
			"action_type": "continuous"
		},
		}

func set_action(action) -> void:
	move_action = clamp(action["move_action"][0], -1.0, 1.0)
\end{lstlisting}

Теперь мы определили наблюдение агента --- положение и скорость мяча в его локальной системе координат. Пространство действий тоже задано: одно непрерывное значение в диапазоне от -1 до +1.

Следующий шаг --- обновить скрипт Player, чтобы использовать действия из AIController. Нажмите значок скрипта рядом с узлом Player и замените содержимое файла \texttt{Player.gd} следующим кодом:

\begin{lstlisting}
extends Node3D

@export var rotation_speed = 3.0
@onready var ball = get_node("../Ball")
@onready var ai_controller = $AIController3D

func _ready():
	ai_controller.init(self)

func game_over():
	ai_controller.done = true
	ai_controller.needs_reset = true

func _physics_process(delta):
	if ai_controller.needs_reset:
		ai_controller.reset()
		ball.reset()
		return

	var movement : float
	if ai_controller.heuristic == "human":
		movement = Input.get_axis("rotate_anticlockwise", "rotate_clockwise")
	else:
		movement = ai_controller.move_action
	rotate_y(movement*delta*rotation_speed)

func _on_area_3d_body_entered(body):
	ai_controller.reward += 1.0
\end{lstlisting}

Теперь необходимо синхронизировать игру, запущенную в Godot, и нейросеть, обучаемую в Python. Для этого Godot RL Agents предоставляет специальный узел. Откройте сцену train.tscn, щёлкните правой кнопкой мыши по корневому узлу и выберите ``Add Child Node. Найдите ``sync и добавьте узел синхронизации Godot RL Agents Sync. Он обеспечивает обмен данными между Python и Godot по TCP.

Обучение можно запускать прямо в редакторе, предварительно запустив Python-процесс обучения командой \texttt{gdrl}.

В этом простом примере агент осваивает достаточно хорошую политику за несколько минут. Чтобы ускорить обучение, выберите узел Sync в сцене train: в инспекторе доступен параметр ``Speed Up:

\CourseFigure{godot6.png}{Узел синхронизации Godot RL Agents: параметр Speed Up --- источник: \protect\href{https://huggingface.co/datasets/huggingface-deep-rl-course/course-images/resolve/main/en/unit9/godot6.png}{[оригинал изображения]}}

Попробуйте установить значение этого параметра равным 8. Это ускорит обучение и особенно пригодится в сложных средах, например в многопользовательском шутере от первого лица, который мы рассмотрим в следующей главе.

\subsubsection{Экспорт модели}\label{export-model}

\href{https://github.com/edbeeching/godot_rl_agents/tree/main?tab=readme-ov-file\#exporting-and-loading-your-trained-agent-in-onnx-format}{Справочная документация}

На время отложим редактор Godot. Чтобы сохранить обученную модель, понадобится выполнить несколько команд в терминале.

В версии библиотеки Godot RL, описанной в исходном материале, доступна экспериментальная поддержка ONNX-моделей для Stable Baselines 3, RLlib и CleanRL.

Например, выберем Stable Baselines 3. Обучите агента при помощи \href{https://github.com/edbeeching/godot_rl_agents/blob/main/examples/stable_baselines3_example.py}{примера SB3} (\href{https://github.com/edbeeching/godot_rl_agents/blob/main/docs/ADV_STABLE_BASELINES_3.md\#train-a-model-from-scratch}{инструкция к скрипту}), включив параметр \texttt{-\/-onnx\_export\_path=model.onnx}.

Пример команды запуска:

\begin{lstlisting}
cd <....> # go into this Godot project directory
python stable_baselines3_example.py --timesteps=100_000 --onnx_export_path=model.onnx --save_model_path=model.zip --save_checkpoint_frequency=20_000 --experiment_name=exp1
\end{lstlisting}

Если всё настроено правильно, в терминале появятся примерно следующие сообщения:

\begin{lstlisting}
No game binary has been provided, please press PLAY in the Godot editor
waiting for remote GODOT connection on port 11008
\end{lstlisting}

\begin{quote}\raggedright
Если при запуске stable\_baselines3\_example возникает ошибка импорта ``ImportError: cannot import name `export\_model\_as\_onnx\' from `godot\_rl.wrappers.onnx.stable\_baselines\_export\'\,, посмотрите решение в \href{https://github.com/edbeeching/godot_rl_agents/issues/203}{этом обсуждении}.
\end{quote}

Теперь вернитесь в редактор Godot и нажмите PLAY в правом верхнем углу. Откроется игровая сцена с обучающимся агентом, а в терминале начнут появляться метрики. Дождитесь завершения обучения. Если всё прошло успешно, в каталоге проекта Godot должен появиться файл \texttt{model.onnx}.

\subsubsection{Использование обученной модели в игре}\label{apply-ai-in-the-game}

Теперь подключим обученную модель к игре!

В редакторе Godot найдите узел Sync в файле \texttt{train.tscn}:

\begin{itemize}
\item
  В выпадающем списке переключите режим управления на \texttt{Onnx\ Inference}.
\item
  Укажите в параметре \texttt{Onnx\ Model\ Path} имя файла модели, в нашем случае \texttt{model.onnx}.
\end{itemize}

Для запуска игры в этом режиме нужна Mono-версия редактора Godot (вариант с поддержкой .NET), которую можно скачать с официального сайта. Также требуется установить \href{https://dotnet.microsoft.com/en-us/download}{.NET}.

С большой вероятностью первая попытка завершится ошибкой. Ниже перечислены возможные проблемы и способы их устранения:

\begin{enumerate}
\def\labelenumi{\arabic{enumi}.}
\item
  Ошибка \texttt{Invalid\ Call.\ Nonexistent\ function\ \textquotesingle{}new\textquotesingle{}\ in\ base\ \textquotesingle{}CSharpScript\textquotesingle{}}: \href{https://github.com/edbeeching/godot_rl_agents/blob/main/docs/TROUBLESHOOTING.md}{решение}.
\item
  Ошибки, связанные с \texttt{onnxruntime} на macOS: \href{https://github.com/microsoft/onnxruntime/issues/9707}{решение}.
\end{enumerate}

\subsubsection{Что ещё можно сделать}\label{theres-more}

Мы рассмотрели лишь малую часть возможностей Godot RL Agents. В библиотеке есть дополнительные пользовательские датчики и камеры, расширяющие информацию, которую получает агент. Посмотрите \href{https://github.com/edbeeching/godot_rl_agents_examples}{примеры}, чтобы узнать больше!

Чтобы узнать об экспорте обученной модели в .onnx для запуска инференса непосредственно в Godot без Python-сервера, а также о других полезных параметрах обучения, изучите \href{https://github.com/edbeeching/godot_rl_agents/blob/main/docs/ADV_STABLE_BASELINES_3.md}{расширенное руководство по SB3}.

\subsection{Автор раздела}\label{author-6}

Автор раздела: \href{https://twitter.com/edwardbeeching}{Эдвард Бичинг (Edward Beeching)}.

\section{Введение в документирование RL-систем}\label{brief-introduction-to-rl-documentation}

В этой дополнительной теме мы рассмотрим вопрос: \textbf{как отслеживать работу и документировать поведение мощных RL-агентов, которые мы обучаем в реальном мире и которые взаимодействуют с людьми?}

По мере того как системы машинного обучения всё сильнее влияют на современную жизнь, \textbf{растёт потребность в их документировании}.

Такая документация может описывать используемые обучающие данные --- где они хранятся, когда были собраны, кто участвовал в сборе и т. д.; либо процесс оптимизации модели --- архитектуру, метрики оценки, связанные статьи и многое другое.

Сегодня всё шире распространяются карточки моделей и паспорта наборов данных. Например, они доступны на Hub (см. \href{https://huggingface.co/docs/hub/model-cards}{документацию}).

Открыв страницу \href{https://huggingface.co/models}{популярной модели на Hub}, можно узнать о процессе её создания.

Такие журналы, относящиеся к конкретной модели или набору данных, обычно заполняют во время их создания, но впоследствии не обновляют, когда модель становится частью развивающейся системы.

\subsection{Зачем нужны отчёты о вознаграждениях (Reward Reports)}\label{motivating-reward-reports}

Системы RL по своей природе оптимизируются по сигналу вознаграждения во времени. Хотя понятие функции вознаграждения можно связать с хорошо изученными областями обучения с учителем (через функцию потерь), наше понимание того, как поведение систем машинного обучения изменяется со временем, пока ограничено.

Поэтому авторы предлагают \href{https://www.notion.so/Brief-introduction-to-RL-documentation-b8cbda5a6f5242338e0756e6bef72af4}{\emph{Reward Reports for Reinforcement Learning --- отчёты о вознаграждениях в RL}}. Название намеренно перекликается с известными работами \emph{Model Cards for Model Reporting} и \emph{Datasheets for Datasets}. Цель --- предложить тип документации, ориентированный на \textbf{человеческие аспекты определения вознаграждения} и \textbf{системы обратной связи, изменяющиеся во времени}.

Опираясь на форматы документации \href{https://arxiv.org/abs/1810.03993}{карточек моделей} и \href{https://arxiv.org/abs/1803.09010}{паспортов наборов данных}, предложенные Митчелл и Гебру с соавторами, мы обосновываем необходимость отчётов Reward Reports для систем ИИ.

\textbf{Reward Reports} --- постоянно обновляемые документы для предполагаемых внедрений RL, фиксирующие выбранные проектные решения.

Однако остаётся много вопросов о применимости такого подхода к разным RL-задачам, трудностях интерпретируемости систем и сходстве уже внедрённых систем обучения с учителем с последовательным принятием решений в RL.

Как минимум Reward Reports дают практикам RL возможность обсудить эти вопросы и приступить к поиску способов их решения.

\subsection{Документирование изменений поведения во времени}\label{capturing-temporal-behavior-with-documentation}

Ключевой элемент документации для RL и ML-систем с обратной связью --- \emph{журнал изменений (change-log)}. Он фиксирует обновления со стороны разработчика (изменившиеся параметры обучения, данные и т. п.) и замеченные пользователями изменения (вредное поведение, неожиданные ответы и т. п.).

В дополнение к журналу предусмотрены события, которые должны запускать обновление отчёта и побуждают наблюдать за последствиями изменений.

\subsection{Участие в проекте}\label{contributing}

Некоторые наиболее влиятельные RL-системы затрагивают множество заинтересованных сторон, но разрабатываются закрыто внутри частных компаний. Поскольку такие компании в значительной степени действуют без специального регулирования, ответственность за документирование часто ложится на общественность.

Если вы хотите принять участие, мы создаём публичные Reward Reports для популярных ML-систем на \href{https://github.com/RewardReports/reward-reports}{GitHub}.

Для дальнейшего чтения обратитесь к \href{https://arxiv.org/abs/2204.10817}{статье о Reward Reports} или посмотрите \href{https://github.com/RewardReports/reward-reports/tree/main/examples}{пример отчёта}.

\subsection{Автор раздела}\label{author-7}

Автор раздела: \href{https://twitter.com/natolambert}{Нейтан Ламберт (Nathan Lambert)}.

\section{Интересные среды для экспериментов}\label{interesting-environments-to-try}

Здесь мы собрали интересные среды, в которых вы можете попробовать обучить своих агентов:

\subsection{DIAMBRA Arena}\label{diambra-arena}

\CourseFigure{diambraarena.png}{DIAMBRA Arena: аркадные среды для экспериментов с RL --- источник: \protect\href{https://huggingface.co/datasets/huggingface-deep-rl-course/course-images/resolve/main/en/unit12/diambraarena.png}{[оригинал изображения]}}

DIAMBRA Arena --- программный пакет с набором качественных сред для исследований и экспериментов в области обучения с подкреплением. Он предоставляет стандартный интерфейс к популярным аркадным играм в эмуляторе, а его Python API полностью соответствует формату OpenAI Gym/Gymnasium, поэтому начать работу с ним легко.

Поддерживаются основные операционные системы (Linux, Windows, macOS); пакет легко установить через \href{https://pypi.org/project/diambra-arena/}{pip}. Использование бесплатно, нужно лишь зарегистрироваться на \href{https://diambra.ai/register/}{официальном сайте}.

Кроме того, в \href{https://github.com/diambra/}{репозитории GitHub} есть примеры основных сценариев использования, которые можно запустить за несколько шагов.

\paragraph{Основные возможности}\label{main-features}

Все среды представляют собой эпизодические RL-задачи с дискретными действиями (кнопки геймпада) и наблюдениями в виде пикселей экрана, дополненными числовыми данными (значения RAM, например полосы здоровья или расположение персонажа на арене).

Поддерживаются одиночный режим (1P) и режим двух игроков (2P). Поэтому платформу удобно использовать для обычного RL, соревновательного многоагентного обучения, соревнований человека и агента, Self-Play, имитационного обучения и обучения с человеком в контуре (Human-in-the-Loop).

\href{https://docs.diambra.ai/envs/games/}{Доступные игры} выбраны из числа популярных ретро-файтингов. Хотя их базовая механика схожа, каждая предлагает свои трудности: разный набор и число персонажей, комбинации приёмов, восстановление здоровья и т. д.

DIAMBRA Arena разработана с учётом совместимости с основными RL-библиотеками. Она предоставляет готовые интерфейсы для \href{https://stable-baselines3.readthedocs.io/en/master/}{Stable Baselines 3} и \href{https://docs.ray.io/en/latest/rllib/index.html}{Ray RLlib}; также доступна устаревшая поддержка Stable Baselines. Примеры приведены в \href{https://docs.diambra.ai/}{официальной документации} и \href{https://github.com/diambra/agents}{репозитории DIAMBRA Agents}. По аналогии можно подключить и другие библиотеки.

\subsubsection{Соревновательная платформа}\label{competition-platform}

DIAMBRA также предоставляет соревновательную платформу, полностью интегрированную с Hugging Face Hub: вы можете публиковать обученных агентов и соревноваться с другими разработчиками со всего мира в турнирах по видеоиграм!

На платформе есть публичная таблица лидеров, в которой пользователи ранжируются по лучшим результатам их агентов в разных средах.

В зависимости от результатов агента можно получать различные достижения.

Поступившие агенты оцениваются, а записи их эпизодов транслируются на \href{https://www.twitch.tv/diambra_ai}{канале DIAMBRA в Twitch}.

\paragraph{Ссылки}\label{references}

Чтобы начать работу с этой средой, воспользуйтесь следующими ресурсами:

\begin{itemize}
\item
  \href{https://docs.diambra.ai/}{Официальная документация}
\item
  \href{https://diambra.ai}{Соревновательная платформа}
\item
  \href{https://github.com/diambra/}{GitHub}
\item
  \href{https://diambra.ai/discord}{Discord}
\end{itemize}

\subsection{MineRL}\label{minerl}

\CourseFigure{minerl.jpg}{MineRL: задачи обучения с подкреплением на основе Minecraft --- источник: \protect\href{https://huggingface.co/datasets/huggingface-deep-rl-course/course-images/resolve/main/en/unit12/minerl.jpg}{[оригинал изображения]}}

MineRL --- библиотека Python с интерфейсом Gym для взаимодействия с игрой Minecraft и наборами данных, собранными из действий людей. Каждый год с использованием библиотеки проводятся соревнования. Подробнее --- на \href{https://minerl.io/}{сайте проекта}.

Чтобы начать работу с этой средой, воспользуйтесь следующими ресурсами:

\begin{itemize}
\item
  \href{https://www.youtube.com/watch?v=z6PTrGifupU}{Что такое MineRL?}
\item
  \href{https://www.youtube.com/watch?v=8yIrWcyWGek}{Первые шаги в MineRL}
\item
  \href{https://minerl.readthedocs.io/en/latest/}{Документация и руководства MineRL}
\end{itemize}

\subsection{Симулятор DonkeyCar}\label{donkeycar-simulator}

\CourseFigure{donkeycar.jpg}{Симуляционная среда DonkeyCar --- источник: \protect\href{https://huggingface.co/datasets/huggingface-deep-rl-course/course-images/resolve/main/en/unit12/donkeycar.jpg}{[оригинал изображения]}}

DonkeyCar --- платформа автономного вождения для любительских радиоуправляемых автомобилей. Эта версия симулятора построена на Unity: используются физика и графика движка, а связь с процессом Python позволяет обученной модели управлять виртуальным автомобилем Donkey.

Чтобы начать работу с этой средой, воспользуйтесь следующими ресурсами:

\begin{itemize}
\item
  \href{https://docs.donkeycar.com/guide/deep_learning/simulator/}{Документация симулятора DonkeyCar}
\item
  \href{https://www.youtube.com/watch?v=ngK33h00iBE}{Плавное вождение: руководство Антонена Раффена, часть 1}
\item
  \href{https://www.youtube.com/watch?v=DUqssFvcSOY}{Плавное вождение: руководство Антонена Раффена, часть 2}
\item
  \href{https://www.youtube.com/watch?v=v8j2bpcE4Rg}{Плавное вождение: руководство Антонена Раффена, часть 3}
\item
  Предварительно обученные агенты:

  \begin{itemize}
  \tightlist
  \item
    \url{https://huggingface.co/araffin/tqc-donkey-mountain-track-v0}
  \item
    \url{https://huggingface.co/araffin/tqc-donkey-avc-sparkfun-v0}
  \item
    \url{https://huggingface.co/araffin/tqc-donkey-minimonaco-track-v0}
  \end{itemize}
\end{itemize}

\subsection{StarCraft II}\label{starcraft-ii}

\CourseFigure{alphastar.jpg}{AlphaStar и StarCraft II как среда для исследования RL --- источник: \protect\href{https://huggingface.co/datasets/huggingface-deep-rl-course/course-images/resolve/main/en/unit12/alphastar.jpg}{[оригинал изображения]}}

StarCraft II --- известная \emph{стратегия в реальном времени}. Компания DeepMind использовала эту игру для исследований в области глубокого RL, создав \href{https://www.deepmind.com/blog/alphastar-mastering-the-real-time-strategy-game-starcraft-ii}{AlphaStar}.

Чтобы начать работу с этой средой, воспользуйтесь следующими ресурсами:

\begin{itemize}
\item
  \href{http://starcraftgym.com/}{StarCraft Gym}
\item
  \href{https://www.youtube.com/watch?v=q59wap1ELQ4}{ИИ учится играть в StarCraft II: руководство по RL}
\end{itemize}

\subsection{Автор раздела}\label{author-8}

Автор раздела: \href{https://twitter.com/ThomasSimonini}{Томас Симонини (Thomas Simonini)}.

\section{Студенческие проекты}\label{student-works}

С момента запуска курса Deep Reinforcement Learning \textbf{многие студенты создали замечательные проекты, с которыми стоит познакомиться и в которых можно принять участие}.

Если вы разработали интересный проект, можете \href{https://github.com/huggingface/deep-rl-class}{добавить его в этот список, создав pull request в репозитории GitHub}.

Проекты \textbf{отсортированы по датам их публикации на этой странице}.

\subsection{Space Scavanger AI}\label{space-scavanger-ai}

Этот проект представляет собой космическую игровую среду с нейросетевым ИИ.

ИИ обучается с помощью методов RL на основе фреймворков Unity ML-Agents и RLlib.

\CourseFigure{spacescavangerai.png}{Студенческий проект Space Scavanger AI --- источник: \protect\href{https://huggingface.co/datasets/huggingface-deep-rl-course/course-images/resolve/main/en/unit12/spacescavangerai.png}{[оригинал изображения]}}

Играть можно здесь: \url{https://swingshuffle.itch.io/spacescalvagerai}

Посмотреть проект Unity: \url{https://github.com/HighExecutor/SpaceScalvagerAI}

\subsection{Neural Nitro}\label{neural-nitro}

\CourseFigure{neuralnitro.png}{Студенческий проект Neural Nitro: гоночная среда --- источник: \protect\href{https://huggingface.co/datasets/huggingface-deep-rl-course/course-images/resolve/main/en/unit12/neuralnitro.png}{[оригинал изображения]}}

В этом проекте Sookeyy создал низкополигональную гоночную игру и научил автомобиль ездить.

Посмотреть демонстрацию: \url{https://sookeyy.itch.io/neuralnitro}

\subsection{Space War}\label{space-war}

\CourseFigure{spacewar.jpg}{Студенческий проект Space War --- источник: \protect\href{https://huggingface.co/datasets/huggingface-deep-rl-course/course-images/resolve/main/en/unit12/spacewar.jpg}{[оригинал изображения]}}

В этом проекте Эрик Донг (Eric Dong) воссоздаёт в Pygame версию игры Space War Билла Сайлера (1985 год) и использует обучение с подкреплением для тренировки ИИ-агентов.

Проект пока находится в разработке!

\subsubsection{Демонстрация}\label{demo}

Версия для разработчиков (Dev/Edge):

\begin{itemize}
\tightlist
\item
  \url{https://e-dong.itch.io/spacewar-dev}
\end{itemize}

Стабильная версия:

\begin{itemize}
\item
  \url{https://e-dong.itch.io/spacewar}
\item
  \url{https://huggingface.co/spaces/EricofRL/SpaceWarRL}
\end{itemize}

\subsubsection{Публикации сообщества}\label{community-blog-posts}

Будет объявлено позже.

\subsubsection{Другие ссылки}\label{other-links}

Исходный код: \url{https://github.com/e-dong/space-war-rl}\\
Блог автора: \url{https://dev.to/edong/space-war-rl-0-series-introduction-25dh}

\subsection{Трансформеры принятия решений для торговли}\label{decision-transformers-for-trading}

В этом проекте студент исследовал обучение Decision Transformer для биржевой торговли. На первом этапе реализовано офлайн-обучение; в следующей версии планируется добавить онлайн-дообучение.

\CourseFigure{decision-transformers-trading.png}{Студенческий проект Decision Transformers for Trading --- источник: \protect\href{https://www.youtube.com/watch?v=w4Bw8WYL8Ps}{Stanford CS25: Decision Transformer и моделирование последовательностей} \protect\href{https://github.com/user-attachments/assets/41e3bc10-2594-4c05-b48f-17666cc8fb1c}{[оригинал изображения]}}

Исходный код: \url{https://github.com/ra9hur/Decision-Transformers-For-Trading}

\end{document}
